\chapter{Random Walks and Electrical Networks}

A token sits on a vertex. At each tick it picks a uniformly random edge out of where it stands and slides along it. That is the whole story of a simple random walk on a graph, and yet it secretly computes a remarkable amount: how long until the token visits a target vertex, how long until it has seen every vertex, how much ``flow'' a graph can route between two points, and which clusters the graph splits into. The reason a one-line stochastic process touches all of these is that it is the same object, viewed three ways: as a Markov chain, as the iterated multiplication of a matrix built from the Laplacian, and --- the surprise that organizes this chapter --- as the diffusion of electric current through a resistor network.

The dictionary is exact. Voltages are harmonic functions and harmonic functions are hitting probabilities. Effective resistance between two vertices is, up to a clean factor of $2m$, the \emph{commute time} of the walk. The energy a unit current dissipates is the reciprocal of conductance, and conductance lower-bounds how fast the walk mixes. Once you have the electrical picture, theorems that look like they need delicate probabilistic arguments fall out of Ohm's law and the fact that current minimizes energy. This chapter builds that dictionary from scratch, proves the central identities (commute time $=2m\,R_{\mathrm{eff}}$, Foster's theorem, the hitting-time formula via the Laplacian pseudoinverse), and traces every one of them by hand on a single small network.

\begin{intuitionbox}
Replace each edge of the graph by a $1$-ohm resistor and wire a $1$-volt battery between two vertices $s$ and $t$. Current floods through the network, splitting at each junction. Two facts make this physical picture a computational engine:
\begin{itemize}
  \item The \emph{voltage} at a vertex is the probability that a random walker started there reaches $s$ before $t$ (after rescaling). Voltages are \emph{harmonic}: each interior vertex's voltage is the average of its neighbors', which is precisely the one-step equation of the walk.
  \item The \emph{effective resistance} $R_{\mathrm{eff}}(s,t)$ --- the voltage you must apply to push one unit of current --- equals $\frac{1}{2m}$ times the expected time for the walk to go from $s$ to $t$ \emph{and back}.
\end{itemize}
So ``how well-connected are $s$ and $t$?'' has one answer wearing two costumes: low resistance (physics) and fast commute (probability).
\end{intuitionbox}

\chapterline

\section{The simple random walk and its transition matrix}

Fix a connected undirected graph $\Gph=(V,E)$ with $n=\abs{V}$ vertices and $m=\abs{E}$ edges. For most of the chapter we take edges unweighted; the weighted case (conductances $w_{ij}>0$) is a drop-in generalization that we flag as we go. Write $\deg(i)$ for the degree of vertex $i$, $\Adjm$ for the adjacency matrix, $\Adjm_D=\mathrm{diag}(\deg(1),\dots,\deg(n))$ for the degree matrix, and $\Lap=\Adjm_D-\Adjm$ for the graph Laplacian introduced in the Laplacian chapter.

\begin{definition}[Simple random walk]\label{ch35:def-srw}
The \emph{simple random walk} on $\Gph$ is the Markov chain $(X_0,X_1,X_2,\dots)$ on $V$ that, from vertex $i$, moves to a uniformly chosen neighbor:
\[
  \Pr[X_{t+1}=j \mid X_t=i] \;=\;
  \begin{cases} \dfrac{1}{\deg(i)} & \{i,j\}\in E,\\[4pt] 0 & \text{otherwise.}\end{cases}
\]
The \emph{transition matrix} is $\Adjm_P \in \R^{n\times n}$ with $(\Adjm_P)_{ij}=\Pr[X_{t+1}=j\mid X_t=i]$. Compactly, $\Adjm_P = \Adjm_D^{-1}\Adjm$.
\end{definition}

Each row of $\Adjm_P$ sums to $1$ (the walker must go somewhere): $\Adjm_P$ is \emph{row-stochastic}. If $p_t \in \R^{n}$ is the row vector of probabilities $\Pr[X_t=i]$, then one step is a vector--matrix product,
\[
  p_{t+1} \;=\; p_t \,\Adjm_P, \qquad\text{so}\qquad p_t \;=\; p_0\,\Adjm_P^{\,t}.
\]
The entire evolution of the walk is the powers of one matrix. This is why so much about random walks reduces to the spectrum of $\Adjm_P$, and why the Laplacian --- so close to $\Adjm_P$ --- is never far away.

\begin{example}[Our running network]\label{ch35:ex-running}
Throughout the chapter we trace a five-vertex graph: a triangle $1\!-\!2\!-\!3$ with a path $3\!-\!4\!-\!5$ hanging off it.

\begin{figure}[H]
\centering
\begin{tikzpicture}[every node/.style={circle,draw,fill=nodeblue!15,minimum size=9mm,inner sep=1pt,font=\small},thick]
  \node (1) at (0,0) {1};
  \node (2) at (1.6,1.1) {2};
  \node (3) at (1.6,-1.1) {3};
  \node (4) at (3.6,-1.1) {4};
  \node (5) at (5.6,-1.1) {5};
  \draw (1)--(2);
  \draw (1)--(3);
  \draw (2)--(3);
  \draw (3)--(4);
  \draw (4)--(5);
\end{tikzpicture}
\caption{Running graph $\Gph$: a triangle $\{1,2,3\}$ with a tail $3\!-\!4\!-\!5$. Degrees: $\deg(1)=\deg(2)=2$, $\deg(3)=3$, $\deg(4)=2$, $\deg(5)=1$. Edges $m=5$.}
\label{ch35:fig-running}
\end{figure}

Reading the edges $\{1,2\},\{1,3\},\{2,3\},\{3,4\},\{4,5\}$:
\[
  \Adjm=\begin{pmatrix}0&1&1&0&0\\1&0&1&0&0\\1&1&0&1&0\\0&0&1&0&1\\0&0&0&1&0\end{pmatrix},
  \quad
  \Adjm_D=\mathrm{diag}(2,2,3,2,1),
  \quad
  \Adjm_P=\begin{pmatrix}
  0&\tfrac12&\tfrac12&0&0\\
  \tfrac12&0&\tfrac12&0&0\\
  \tfrac13&\tfrac13&0&\tfrac13&0\\
  0&0&\tfrac12&0&\tfrac12\\
  0&0&0&1&0
  \end{pmatrix}.
\]
Vertex $5$ is a dead end: whenever the walk reaches it, the next step is forced back to $4$. Keep $\Adjm_P$ in view; we will find its stationary distribution, then its hitting and commute times, all by hand.
\end{example}

\subsection{The stationary distribution}

\begin{definition}[Stationary distribution]\label{ch35:def-stationary}
A probability row vector $\pi$ is \emph{stationary} for the walk if $\pi\,\Adjm_P=\pi$. It is a fixed point: start the walk distributed as $\pi$ and it stays distributed as $\pi$ forever.
\end{definition}

For the simple random walk the stationary distribution is not uniform; it is proportional to degree. Hubs get visited more because more edges point at them.

\begin{theorem}[Degree-proportional stationarity]\label{ch35:thm-stationary}
On a connected graph the vector
\[
  \pi_i \;=\; \frac{\deg(i)}{2m}
\]
is stationary for the simple random walk, and it is the unique stationary distribution.
\end{theorem}

\begin{proof}
Compute the $j$-th entry of $\pi\Adjm_P$:
\[
  (\pi\Adjm_P)_j \;=\; \sum_{i}\pi_i (\Adjm_P)_{ij}
  \;=\; \sum_{i:\{i,j\}\in E}\frac{\deg(i)}{2m}\cdot\frac{1}{\deg(i)}
  \;=\; \sum_{i:\{i,j\}\in E}\frac{1}{2m}
  \;=\;\frac{\deg(j)}{2m}\;=\;\pi_j.
\]
The two $\deg(i)$ factors cancel --- the engine of the whole calculation --- leaving one $\frac{1}{2m}$ per neighbor of $j$, of which there are $\deg(j)$. So $\pi\Adjm_P=\pi$. Uniqueness: on a connected graph the chain is irreducible, and the Perron--Frobenius theorem for stochastic matrices gives a unique stationary distribution. (If the graph is also non-bipartite, the chain is aperiodic and $p_t\to\pi$ from every start; we return to convergence in Section~\ref{ch35:mixing}.)
\end{proof}

This is a special case of \emph{reversibility}. Define the \emph{detailed-balance} condition $\pi_i (\Adjm_P)_{ij}=\pi_j(\Adjm_P)_{ji}$. With $\pi_i=\deg(i)/2m$ both sides equal $\frac{1}{2m}\mathbf{1}[\{i,j\}\in E]$, so the simple random walk is reversible. Reversibility is exactly what makes the electrical analogy work: a reversible chain is a resistor network in disguise.

\begin{example}[Stationary distribution of $\Gph$]\label{ch35:ex-stationary}
With $2m=10$, the stationary distribution of our running graph is
\[
  \pi=\Big(\tfrac{2}{10},\tfrac{2}{10},\tfrac{3}{10},\tfrac{2}{10},\tfrac{1}{10}\Big)
  =(0.2,\,0.2,\,0.3,\,0.2,\,0.1).
\]
The degree-$3$ hub vertex $3$ is visited most; the leaf vertex $5$ least. Sanity check one coordinate: $(\pi\Adjm_P)_5 = \pi_4\cdot(\Adjm_P)_{45}=0.2\cdot\tfrac12=0.1=\pi_5$. Good.
\end{example}

\begin{trapbox}
Two failure modes break ``$\pi_i=\deg(i)/2m$ and the walk converges to it.''
\emph{Disconnected graph}: stationarity is no longer unique --- each component carries its own $\pi$. \emph{Bipartite graph}: the chain is periodic (period $2$); $\pi$ is still stationary, but $p_t$ oscillates between the two sides and never settles. A single odd cycle, like the triangle in our running graph, kills bipartiteness and restores convergence. The \emph{lazy} walk $\Adjm_P'=\tfrac12(\Adjm_I+\Adjm_P)$ --- flip a coin, stay put on heads --- is aperiodic on \emph{any} graph and has the same $\pi$; it is the standard fix.
\end{trapbox}

\section{Hitting times and cover times}

\begin{definition}[Hitting and commute times]\label{ch35:def-hitting}
For vertices $a,b$, let $T_b=\min\{t\ge 0: X_t=b\}$ be the first time the walk reaches $b$. The \emph{hitting time} is the expected first arrival starting from $a$,
\[
  H(a,b)\;=\;\mathbb{E}\big[T_b \mid X_0=a\big].
\]
The \emph{commute time} is the round trip, $C(a,b)=H(a,b)+H(b,a)$. The \emph{cover time} from $a$ is $\mathbb{E}[\max_{v}T_v\mid X_0=a]$, the expected time to visit \emph{every} vertex; the cover time of $\Gph$ is the maximum over start vertices.
\end{definition}

Hitting time is generally \emph{not} symmetric: $H(a,b)\ne H(b,a)$. In our running graph, getting from the leaf $5$ into the triangle is easy, but wandering out to the leaf $5$ from the triangle is slow, because the walk keeps getting pulled back into the dense core. The commute time, being a round trip, is symmetric by construction --- and that symmetry is the first hint of the resistance connection.

\subsection{Hitting times by first-step analysis}

The workhorse identity is \emph{first-step analysis}: condition on the first move. From $a\ne b$ the walk takes one step (cost $1$) to a uniform neighbor, then must still reach $b$:
\begin{equation}\label{ch35:eq-firststep}
  H(a,b) \;=\; 1+\frac{1}{\deg(a)}\sum_{c:\{a,c\}\in E} H(c,b),
  \qquad H(b,b)=0.
\end{equation}
Fixing the target $b$, this is a linear system in the $n-1$ unknowns $\{H(a,b):a\ne b\}$. Writing $h$ for the vector of hitting times to $b$ and deleting the row/column for $b$, \eqref{ch35:eq-firststep} reads $(\Adjm_I-\Adjm_P)\,h=\mathbf{1}$ on the non-$b$ coordinates, a Laplacian-type system we solve below in closed form.

\begin{example}[Hitting times on the tail of $\Gph$]\label{ch35:ex-hitting}
Let us compute hitting times to vertex $5$ by hand. By symmetry of the triangle, vertices $1$ and $2$ are interchangeable with respect to the target $5$, so $H(1,5)=H(2,5)$. Write $h_a=H(a,5)$. The equations \eqref{ch35:eq-firststep} are:
\begin{align*}
  h_1 &= 1+\tfrac12 h_2+\tfrac12 h_3,\\
  h_2 &= 1+\tfrac12 h_1+\tfrac12 h_3,\\
  h_3 &= 1+\tfrac13 h_1+\tfrac13 h_2+\tfrac13 h_4,\\
  h_4 &= 1+\tfrac12 h_3+\tfrac12\cdot 0.
\end{align*}
Using $h_1=h_2$, the first equation gives $h_1=1+\tfrac12 h_1+\tfrac12 h_3$, so $h_1=2+h_3$. Substituting into the third, $h_3=1+\tfrac23 h_1+\tfrac13 h_4=1+\tfrac23(2+h_3)+\tfrac13 h_4$, i.e. $\tfrac13 h_3=\tfrac73+\tfrac13 h_4$, so $h_3=7+h_4$. The fourth gives $h_4=1+\tfrac12 h_3$. Combine: $h_4=1+\tfrac12(7+h_4)=\tfrac92+\tfrac12 h_4$, hence $h_4=9$, then $h_3=16$, and $h_1=h_2=18$.

So $H(1,5)=H(2,5)=18$, $H(3,5)=16$, $H(4,5)=9$. The walk needs $18$ steps on average to crawl out to the leaf from a far triangle vertex (Figure~\ref{ch35:fig-hitting}). We will recompute $H(4,5)=9$ instantly from resistance in Example~\ref{ch35:ex-commute}.

\begin{figure}[H]
\centering
\begin{tikzpicture}[every node/.style={circle,draw,fill=nodeblue!15,minimum size=9mm,inner sep=1pt,font=\small},thick]
  \node (1) at (0,0) {1};
  \node (2) at (1.6,1.1) {2};
  \node (3) at (1.6,-1.1) {3};
  \node (4) at (3.6,-1.1) {4};
  \node (5) at (5.6,-1.1) {5};
  \draw (1)--(2); \draw (1)--(3); \draw (2)--(3); \draw (3)--(4); \draw (4)--(5);
  \node[draw=none,fill=none,font=\scriptsize,above left=0pt of 1] {$18$};
  \node[draw=none,fill=none,font=\scriptsize,above=1pt of 2] {$18$};
  \node[draw=none,fill=none,font=\scriptsize,below left=0pt of 3] {$16$};
  \node[draw=none,fill=none,font=\scriptsize,below=1pt of 4] {$9$};
  \node[draw=none,fill=none,font=\scriptsize,below=1pt of 5] {$0$};
\end{tikzpicture}
\caption{Hitting times $H(a,5)$ to the leaf, written beside each vertex. They grow monotonically with distance from $5$ and are equal at the symmetric pair $\{1,2\}$. The target $5$ has $H(5,5)=0$.}
\label{ch35:fig-hitting}
\end{figure}
\end{example}

\begin{labbox}
Compute the hitting times \emph{to} the hub vertex $3$ on the running graph, $H(a,3)$ for $a\in\{1,2,4,5\}$, by first-step analysis. You should find $H(1,3)=H(2,3)$ small and $H(5,3)$ the largest. Then check that $H(3,5)+H(5,3)$ matches the commute time $C(3,5)$ we get from resistance in Section~\ref{ch35:commute}. Doing the two computations independently and watching them agree is the fastest way to believe the commute-resistance theorem.
\end{labbox}

\subsection{Bounds on hitting and cover times}

How large can these times be? For hitting time, a clean universal bound comes straight from commute time and resistance (Section~\ref{ch35:commute}); here we record the headline.

\begin{theorem}[Polynomial hitting time]\label{ch35:thm-hitbound}
On any connected graph with $n$ vertices and $m$ edges, $H(a,b)\le 2m\cdot(n-1) < n^3$ for all $a,b$, and there exist graphs (the ``lollipop'') where $H(a,b)=\Th(n^3)$.
\end{theorem}

\begin{proof}[Proof sketch]
We show in Theorem~\ref{ch35:thm-commute} that $C(a,b)=2m\,R_{\mathrm{eff}}(a,b)$, and effective resistance is at most the graph distance $\le n-1$ (resistors in series). So $H(a,b)\le C(a,b)\le 2m(n-1)\le n(n-1)^2<n^3$. The lollipop graph --- a clique on $n/2$ vertices joined by a path of length $n/2$ --- forces the walk to escape the sticky clique and then random-walk down the path, costing $\Th(n^2)$ to leave times $\Th(n)$ path length, i.e. $\Th(n^3)$; this is the worst case.
\end{proof}

The cover time inherits a logarithmic overhead on top of the worst hitting time.

\begin{theorem}[Matthews bound]\label{ch35:thm-matthews}
Let $H_{\max}=\max_{a,b}H(a,b)$. The cover time satisfies
\[
  \mathrm{cov}(\Gph)\;\le\; H_{\max}\cdot\Big(1+\tfrac12+\cdots+\tfrac1n\Big)\;=\;H_{\max}\,\Oh(\log n).
\]
\end{theorem}

\begin{proof}[Idea]
Order the vertices in a uniformly random order $v_1,\dots,v_n$ and let $L_k$ be the time at which the $k$-th \emph{distinct} ``record'' vertex (the last among $v_1,\dots,v_k$ to be hit) is covered. A coupon-collector argument bounds $\mathbb{E}[\text{cover time}]$ by $\sum_k \frac{1}{k}H_{\max}$, because the chance that $v_k$ is the newest record is $1/k$. Summing the harmonic series gives the $\Oh(\log n)$ factor. The matching lower bound (cover time $\ge H_{\min}\cdot H_n$ with $H_{\min}$ the min over a clever set) is the other half of Matthews' theorem.
\end{proof}

\begin{complexitybox}
For an $n$-vertex, $m$-edge connected graph: hitting time $H(a,b)=\Oh(mn)=\Oh(n^3)$, tight for the lollipop. Cover time $\mathrm{cov}(\Gph)=\Oh(mn\log n)$; it is $\Th(n\log n)$ for the complete graph and $\Th(n^2)$ for the path/cycle. Computing one hitting time exactly costs one Laplacian solve, $\Oh(n^3)$ by Gaussian elimination or $\tilde\Oh(m)$ by the near-linear Laplacian solvers of the Laplacian-solver chapter.
\end{complexitybox}

\section{The electrical-network analogy}\label{ch35:electrical}

Now the change of viewpoint. Turn the graph into a resistor network: each edge $\{i,j\}$ becomes a resistor of resistance $r_{ij}$, equivalently a \emph{conductance} $c_{ij}=1/r_{ij}$. For the simple random walk all conductances are $1$ (unit resistors); a weighted walk with edge weights $w_{ij}$ corresponds to conductances $c_{ij}=w_{ij}$. The vertex conductance is $c_i=\sum_j c_{ij}=\deg(i)$ in the unit case.

\begin{figure}[H]
\centering
\begin{tikzpicture}[font=\small,thick]
    \node[circle,draw,fill=nodeblue!15,minimum size=8mm] (a) at (0,0){$i$};
    \node[circle,draw,fill=nodeblue!15,minimum size=8mm] (b) at (3.5,0){$j$};
    \draw (a) -- (1.2,0) -- (1.4,0.28) -- (1.75,-0.28) -- (2.1,0.28) -- (2.3,0) -- (b);
    \node at (1.75,0.6) {$r_{ij}=1$};
    \node at (-0.1,-0.7) {$v(i)$};
    \node at (3.6,-0.7) {$v(j)$};
    \node at (1.75,-1.35) {current $\;i_{ij}=\dfrac{v(i)-v(j)}{r_{ij}}$ (Ohm)};
\end{tikzpicture}
\caption{One edge as a resistor. Ohm's law: current equals voltage drop over resistance. Kirchhoff's current law: at every interior vertex, current in equals current out.}
\label{ch35:fig-ohm}
\end{figure}

\subsection{Voltages are harmonic functions}

Attach a voltage $v(i)$ to each vertex. \emph{Ohm's law} says the current along edge $\{i,j\}$ from $i$ to $j$ is $i_{ij}=c_{ij}(v(i)-v(j))$. \emph{Kirchhoff's current law} says that at every vertex where no external current is injected, the net current is zero:
\[
  \sum_{j:\{i,j\}\in E} c_{ij}\big(v(i)-v(j)\big)=0
  \quad\Longleftrightarrow\quad
  v(i)=\frac{\sum_j c_{ij}\,v(j)}{\sum_j c_{ij}}
  \;=\;\sum_j (\Adjm_P)_{ij}\,v(j).
\]
The right-hand equation is the definition of a \emph{harmonic} function: the value at $i$ is the (conductance-weighted) average of its neighbors. And $\sum_j(\Adjm_P)_{ij}v(j)$ is exactly one step of the random walk applied to $v$. So:

\begin{whybox}
A function is harmonic at vertex $i$ $\iff$ Kirchhoff's law holds at $i$ $\iff$ $v(i)=(\Adjm_P v)(i)$, i.e. $v$ is unchanged by one step of the walk's averaging. The Laplacian packages all three: $v$ is harmonic at $i$ exactly when $(\Lap v)(i)=0$, since $(\Lap v)(i)=\sum_j c_{ij}(v(i)-v(j))$ is the net current out of $i$. Harmonic functions, zero-net-current vertices, and walk-invariant values are one concept in three languages.
\end{whybox}

Now \emph{inject} current: push a current of $f$ in at vertex $s$ and pull $f$ out at $t$. Then $\Lap v=f(\mathbf{1}_s-\mathbf{1}_t)$, a Laplacian system whose right-hand side is the net injected current. Because $\Lap$ has the all-ones vector in its kernel, voltages are determined only up to an additive constant (ground one vertex), exactly as in physics.

\begin{theorem}[Voltages solve the Dirichlet problem]\label{ch35:thm-dirichlet}
Fix disjoint boundary sets with prescribed voltages (e.g. $v(s)=1$, $v(t)=0$). There is a unique function $v$ that takes those boundary values and is harmonic at every interior vertex. Equivalently, the voltage $v(a)$ equals the probability that a random walk started at $a$ reaches $s$ before $t$.
\end{theorem}

\begin{proof}
\emph{Existence/uniqueness.} Harmonicity at interior vertices plus fixed boundary values is the linear system $(\Lap v)(a)=0$ for interior $a$, with $v(s)=1,v(t)=0$. The interior block of $\Lap$ (rows/cols of interior vertices) is a symmetric \emph{diagonally dominant} matrix that is strictly diagonally dominant in any row touching the boundary; on a connected graph it is invertible, so the system has a unique solution.

\emph{Probabilistic identity.} Let $p(a)=\Pr[\text{reach }s\text{ before }t\mid X_0=a]$. Clearly $p(s)=1$, $p(t)=0$. For interior $a$, first-step analysis gives $p(a)=\sum_j(\Adjm_P)_{aj}p(j)$ --- the walk steps to a random neighbor and continues --- so $p$ is harmonic with the same boundary values as $v$. By uniqueness, $p=v$. The voltage you read off a resistor network \emph{is} a hitting probability.
\end{proof}

This is the maximum principle in action: a harmonic function attains its extremes on the boundary (no interior vertex can exceed all its neighbors, being their average), so the walk-as-voltage is squeezed between $0$ and $1$, as a probability must be.

\subsection{Effective resistance and energy}

\begin{definition}[Effective resistance]\label{ch35:def-reff}
The \emph{effective resistance} $R_{\mathrm{eff}}(s,t)$ between $s$ and $t$ is the voltage drop $v(s)-v(t)$ needed to drive a \emph{unit} current ($f=1$) from $s$ to $t$. Equivalently, $R_{\mathrm{eff}}(s,t)=v(s)-v(t)$ where $\Lap v=\mathbf{1}_s-\mathbf{1}_t$.
\end{definition}

Effective resistance collapses the whole network, however tangled, into a single equivalent resistor between $s$ and $t$. Two familiar rules compute it for simple networks: resistors in \emph{series} add ($R=r_1+r_2$), resistors in \emph{parallel} add reciprocally ($\frac1R=\frac1{r_1}+\frac1{r_2}$). More edges between two points can only lower resistance --- which matches the intuition that more routes make $s$ and $t$ better connected.

There is a clean variational characterization. Recall the quadratic form $f^\top\Lap f=\sum_{\{i,j\}}c_{ij}(f_i-f_j)^2$ from the Laplacian chapter; physically it is the power dissipated. Effective resistance is the minimum energy to route a unit flow:

\begin{theorem}[Thomson's principle]\label{ch35:thm-thomson}
Among all unit $s$--$t$ flows (functions on oriented edges obeying conservation at interior vertices and net $1$ out of $s$, into $t$), the electrical current is the unique minimizer of the dissipated energy $\mathcal{E}(i)=\sum_{e}r_e\,i_e^2$, and the minimum value equals $R_{\mathrm{eff}}(s,t)$.
\end{theorem}

\begin{proof}[Idea]
Energy $\mathcal{E}$ is a strictly convex function on the affine space of unit flows, so it has a unique minimizer. Setting the gradient to zero subject to conservation produces exactly Kirchhoff's voltage law (the flow is a gradient of potentials), which characterizes the electrical current. Evaluating $\mathcal{E}$ at the current, $\sum_e r_e i_e^2=\sum_e i_e\cdot(\text{$v$-drop})=$ (by conservation, telescoping) $1\cdot(v(s)-v(t))=R_{\mathrm{eff}}(s,t)$.
\end{proof}

Thomson's principle has a dual, \emph{Dirichlet's principle}: voltages minimize $f^\top\Lap f$ subject to the boundary conditions. The two are Lagrangian duals, and together they give the monotonicity law (Rayleigh): cutting an edge can only \emph{raise} effective resistance; shorting two vertices can only lower it. We use this to bound resistances by hand.

\begin{example}[Effective resistance in $\Gph$]\label{ch35:ex-reff}
Compute $R_{\mathrm{eff}}(3,5)$ on the running graph. The path $3\!-\!4\!-\!5$ is two unit resistors in series, $r=2$, and it is the \emph{only} route from $3$ to $5$ (vertex $4$ is a cut vertex), so $R_{\mathrm{eff}}(3,5)=2$.

Now $R_{\mathrm{eff}}(1,2)$. There are two routes from $1$ to $2$: the direct edge ($r=1$) and the detour $1\!-\!3\!-\!2$ ($r=2$). The dangling tail $3\!-\!4\!-\!5$ carries no current (it is a dead end with no sink), so it does not affect the resistance. Two parallel resistors, $1$ and $2$:
\[
  \frac{1}{R_{\mathrm{eff}}(1,2)}=\frac11+\frac12=\frac32
  \quad\Longrightarrow\quad
  R_{\mathrm{eff}}(1,2)=\frac23.
\]
Finally $R_{\mathrm{eff}}(1,5)$: from $1$ to $3$ is the parallel combination of edge $1\!-\!3$ ($r=1$) with the path $1\!-\!2\!-\!3$ ($r=2$), giving $\frac23$; then $3\!-\!4\!-\!5$ in series adds $2$. So $R_{\mathrm{eff}}(1,5)=\frac23+2=\frac83$, as Figure~\ref{ch35:fig-reduce} shows step by step.

\begin{figure}[H]
\centering
\begin{tikzpicture}[font=\small,thick,
    res/.style={draw,fill=white,rounded corners=1pt,inner sep=2pt,font=\scriptsize}]
  % Stage 1: original 1--3 part
  \node[circle,draw,fill=nodeblue!15,minimum size=7mm] (a1) at (0,0) {$1$};
  \node[circle,draw,fill=nodeblue!15,minimum size=7mm] (a3) at (2.4,0) {$3$};
  \draw (a1) to[bend left=45] node[res]{$1$} (a3);
  \draw (a1) to[bend right=45] node[res,below]{$2$} (a3);
  \node[font=\scriptsize] at (1.2,-1.25) {edge $\|$ path $1\,2\,3$};
  % arrow
  \node at (3.5,0) {$\Longrightarrow$};
  % Stage 2: reduced 1--3 then 3-4-5
  \node[circle,draw,fill=nodeblue!15,minimum size=7mm] (b1) at (4.6,0) {$1$};
  \node[circle,draw,fill=nodeblue!15,minimum size=7mm] (b3) at (7.0,0) {$3$};
  \node[circle,draw,fill=nodeblue!15,minimum size=7mm] (b5) at (9.4,0) {$5$};
  \draw (b1) -- node[res]{$\tfrac23$} (b3);
  \draw (b3) -- node[res]{$2$} (b5);
  \node[font=\scriptsize] at (8.2,-1.25) {series: $\tfrac23+2=\tfrac83$};
\end{tikzpicture}
\caption{Reducing $R_{\mathrm{eff}}(1,5)$. Left: between $1$ and $3$ the direct edge ($r{=}1$) is in parallel with the path $1\!-\!2\!-\!3$ ($r{=}2$), collapsing to $\tfrac23$. Right: that equivalent resistor is in series with the tail $3\!-\!4\!-\!5$ ($r{=}2$), giving $R_{\mathrm{eff}}(1,5)=\tfrac83$.}
\label{ch35:fig-reduce}
\end{figure}
\end{example}

\subsection{Resistance via the Laplacian pseudoinverse}

For a general network we need a formula, not a picture. Let $\Lap^{+}$ be the Moore--Penrose pseudoinverse of the Laplacian. Solving $\Lap v=\mathbf{1}_s-\mathbf{1}_t$ on the orthogonal complement of $\mathbf{1}$ gives $v=\Lap^{+}(\mathbf{1}_s-\mathbf{1}_t)$, and reading off the drop:

\begin{theorem}[Resistance formula]\label{ch35:thm-reffform}
With $\mathbf{1}_s$ the indicator of $s$,
\[
  R_{\mathrm{eff}}(s,t)
  \;=\;(\mathbf{1}_s-\mathbf{1}_t)^\top \Lap^{+}(\mathbf{1}_s-\mathbf{1}_t)
  \;=\;\Lap^{+}_{ss}+\Lap^{+}_{tt}-2\Lap^{+}_{st}.
\]
Equivalently, $R_{\mathrm{eff}}(s,t)=\sum_{k\ge 2}\dfrac{1}{\lambda_k}\big(u_k(s)-u_k(t)\big)^2$, where $0=\lambda_1<\lambda_2\le\cdots\le\lambda_n$ are the Laplacian eigenvalues and $u_k$ the orthonormal eigenvectors.
\end{theorem}

\begin{proof}
Write $\chi=\mathbf{1}_s-\mathbf{1}_t$, which is orthogonal to $\mathbf{1}$ (entries sum to $0$), so $\Lap^+\Lap\chi=\chi$. Then $v=\Lap^+\chi$ solves $\Lap v=\chi$, and $R_{\mathrm{eff}}=v(s)-v(t)=\chi^\top v=\chi^\top\Lap^+\chi$. Expanding $\chi^\top\Lap^+\chi=(\mathbf{1}_s-\mathbf{1}_t)^\top\Lap^+(\mathbf{1}_s-\mathbf{1}_t)$ gives the entrywise form. The spectral form follows from $\Lap^+=\sum_{k\ge2}\lambda_k^{-1}u_ku_k^\top$ (the eigendecomposition of the Laplacian restricted to $\mathbf{1}^\perp$) and $\chi^\top u_k=u_k(s)-u_k(t)$.
\end{proof}

The spectral form is the bridge to spectral graph theory: resistance is a sum over the Laplacian spectrum, weighted by inverse eigenvalues, so small eigenvalues (poor connectivity, see Cheeger) inflate resistance, exactly as poor connectivity should.

\begin{remark}[Resistance is a metric]\label{ch35:rem-metric}
$R_{\mathrm{eff}}$ is a genuine metric on $V$: symmetric, zero iff $s=t$, and it satisfies the triangle inequality (a one-line consequence of the energy formulation). The map $s\mapsto (\Lap^+)^{1/2}\mathbf{1}_s$ embeds the vertices into Euclidean space so that squared distances are resistances --- the \emph{resistive embedding}, used in graph sparsification and machine learning.
\end{remark}

\section{Commute time equals \texorpdfstring{$2m$}{2m} times resistance}\label{ch35:commute}

Here is the theorem the whole chapter has been pointing at. It welds the probabilistic quantity (commute time) to the electrical one (effective resistance) with a constant that depends only on the number of edges.

\begin{theorem}[Chandra--Raghavan--Ruzzo--Smolensky--Tiwari]\label{ch35:thm-commute}
For every pair $s,t$ in a connected graph,
\[
  C(s,t)\;=\;H(s,t)+H(t,s)\;=\;2m\cdot R_{\mathrm{eff}}(s,t).
\]
For a weighted graph with total conductance $2W=\sum_{i}c_i$, the constant $2m$ becomes $2W$.
\end{theorem}

\begin{proof}
The cleanest route superposes two electrical experiments. Fix a target $t$. Define a function $\phi_a=H(a,t)$, the hitting time to $t$ from $a$. First-step analysis \eqref{ch35:eq-firststep} says, for $a\ne t$,
\[
  \phi_a=1+\sum_j (\Adjm_P)_{aj}\phi_j
  \;\Longleftrightarrow\;
  \sum_j (\Adjm_P)_{aj}(\phi_a-\phi_j)=1
  \;\Longleftrightarrow\;
  \frac{1}{\deg(a)}\sum_{j\sim a}(\phi_a-\phi_j)=1.
\]
Multiply by $\deg(a)$: $\sum_{j\sim a}(\phi_a-\phi_j)=\deg(a)$, i.e. $(\Lap\phi)_a=\deg(a)$ for all $a\ne t$, while $\phi_t=0$. So $\phi$ solves $\Lap\phi=\mathbf{d}-\alpha\,\mathbf{1}_t$ where $\mathbf{d}$ is the degree vector and $\alpha$ absorbs the deficit at $t$. Summing all entries of $\Lap\phi$ (which total $0$, since every column of $\Lap$ sums to $0$) forces $\alpha=\sum_a\deg(a)=2m$. Hence
\[
  \Lap\,\phi \;=\; \mathbf{d}-2m\,\mathbf{1}_t,
  \qquad \phi_a=H(a,t).
\]
Read this electrically: $\phi=H(\cdot,t)$ is the voltage when we inject $\deg(a)$ units of current at every vertex $a$ and extract all $2m$ units at $t$. By symmetry, $\psi_a=H(a,s)$ solves $\Lap\psi=\mathbf{d}-2m\,\mathbf{1}_s$: inject $\deg(a)$ everywhere, extract $2m$ at $s$.

Subtract the two systems: $\Lap(\phi-\psi)=2m(\mathbf{1}_s-\mathbf{1}_t)$. The injected-everywhere currents cancel, leaving a clean $s$--$t$ current of magnitude $2m$. By Definition~\ref{ch35:def-reff} (resistance scales the drop linearly with current), the voltage drop of $\phi-\psi$ between $s$ and $t$ is $2m\,R_{\mathrm{eff}}(s,t)$. Evaluate that drop directly:
\[
  (\phi-\psi)_s-(\phi-\psi)_t
  =\big(\underbrace{H(s,t)}_{\phi_s}-\underbrace{H(s,s)}_{\psi_s=0}\big)
   -\big(\underbrace{H(t,t)}_{\phi_t=0}-\underbrace{H(t,s)}_{\psi_t}\big)
  =H(s,t)+H(t,s)=C(s,t).
\]
Therefore $C(s,t)=2m\,R_{\mathrm{eff}}(s,t)$.
\end{proof}

\begin{aestheticbox}
Look at what just happened. Hitting time is an asymmetric, path-dependent expectation; nothing about its definition is symmetric in $s$ and $t$. Yet its symmetrization, the commute time, equals $2m$ times a quantity (resistance) computed by Ohm's law with no probability in sight. The proof never simulates a walk: it solves two linear systems and subtracts. This is the payoff of reversibility --- the chain \emph{is} a resistor network, so the round trip is forced to be electrical. Every later hitting-time bound in this book routes through this one identity.
\end{aestheticbox}

\begin{example}[Commute times in $\Gph$, finished by hand]\label{ch35:ex-commute}
We computed $R_{\mathrm{eff}}(3,5)=2$ (Example~\ref{ch35:ex-reff}) and $2m=10$, so
\[
  C(3,5)=H(3,5)+H(5,3)=10\cdot 2=20.
\]
We already know $H(3,5)=16$ from Example~\ref{ch35:ex-hitting}, so $H(5,3)=20-16=4$. Getting from the leaf into the hub takes $4$ steps on average; the reverse takes $16$. The asymmetry is stark and the round trip is exactly $2m\,R_{\mathrm{eff}}$, as promised.

Likewise $R_{\mathrm{eff}}(4,5)=1$ (single edge), so $C(4,5)=10$, and since $H(5,4)=1$ (forced step), $H(4,5)=9$ --- matching the brute-force value from Example~\ref{ch35:ex-hitting} with zero extra algebra. Finally $C(1,2)=10\cdot\tfrac23=\tfrac{20}{3}$ and $C(1,5)=10\cdot\tfrac83=\tfrac{80}{3}$.
\end{example}

\subsection{Foster's theorem: a global resistance identity}

Summing resistances over edges gives an identity so clean it feels like a typo.

\begin{theorem}[Foster]\label{ch35:thm-foster}
In any connected graph on $n$ vertices,
\[
  \sum_{\{i,j\}\in E} R_{\mathrm{eff}}(i,j)\;=\;n-1.
\]
\end{theorem}

\begin{proof}
By the resistance formula, $R_{\mathrm{eff}}(i,j)=(\mathbf{1}_i-\mathbf{1}_j)^\top\Lap^+(\mathbf{1}_i-\mathbf{1}_j)$. For an edge, $\mathbf{1}_i-\mathbf{1}_j$ is the column of the signed incidence matrix $\Adjm_B$ for that edge (up to sign), so $R_{\mathrm{eff}}(i,j)=\big(\Adjm_B^\top\Lap^+\Adjm_B\big)_{ee}$. Summing over edges is the trace:
\[
  \sum_e R_{\mathrm{eff}}(e)=\mathrm{tr}\big(\Adjm_B^\top\Lap^+\Adjm_B\big)
  =\mathrm{tr}\big(\Lap^+\Adjm_B\Adjm_B^\top\big)
  =\mathrm{tr}(\Lap^+\Lap)
  =\rank(\Lap)=n-1,
\]
using $\Lap=\Adjm_B\Adjm_B^\top$ (the incidence factorization) and $\mathrm{tr}(\Lap^+\Lap)=\rank\Lap=n-1$ for a connected graph.
\end{proof}

\begin{example}[Foster on $\Gph$]\label{ch35:ex-foster}
The five edges have resistances $R_{\mathrm{eff}}(1,2)=R_{\mathrm{eff}}(1,3)=R_{\mathrm{eff}}(2,3)=\tfrac23$ (each triangle edge sees the other two-edge path in parallel) and $R_{\mathrm{eff}}(3,4)=R_{\mathrm{eff}}(4,5)=1$ (bridges). Sum:
\[
  3\cdot\tfrac23+1+1=2+2=4=n-1.\quad\checkmark
\]
A bridge always has $R_{\mathrm{eff}}=1$ (no parallel path), and indeed the two bridge edges contribute $2$ while the triangle contributes $2$, totaling exactly $n-1=4$.
\end{example}

\begin{connectionbox}
Foster's theorem is the seed of \emph{spectral sparsification}. Spielman and Srivastava showed that sampling each edge with probability proportional to $w_e R_{\mathrm{eff}}(e)$ --- and Foster says these ``effective weights'' sum to $n-1$, so $\Oh(n\log n/\eps^2)$ samples suffice --- produces a sparse graph whose Laplacian approximates the original's quadratic form to within $1\pm\eps$. The same resistances drive the Laplacian-solver and electrical-flow machinery of the solvers chapter. Random walks, resistances, and sparsifiers are one circle of ideas. The probability companion volume's chapter on Markov chains develops the same first-step analysis abstractly; here it wears a graph.
\end{connectionbox}

\section{Computing the quantities: an algorithm}

All the named quantities reduce to Laplacian solves. Below is the pipeline for hitting times, commute times, and resistances; it is the practical recipe behind the theory.

\begin{algorithm}[H]
\DontPrintSemicolon
\caption{Random-walk quantities via Laplacian solves}
\KwIn{Connected graph $\Gph=(V,E)$ with conductances $c_{ij}$; query vertices $s,t$.}
\KwOut{$R_{\mathrm{eff}}(s,t)$, $C(s,t)$, and hitting times $H(\cdot,t)$.}
Assemble $\Lap=\Adjm_D-\Adjm$ (weighted: $\Adjm_D=\mathrm{diag}(c_i)$)\;
$2W \leftarrow \sum_i c_i$ \tcp*{$=2m$ in the unit-resistor case}
\tcp{Effective resistance: one solve}
Solve $\Lap\,v=\mathbf{1}_s-\mathbf{1}_t$ for $v\perp\mathbf{1}$ (ground $v$ at any vertex)\;
$R_{\mathrm{eff}}(s,t)\leftarrow v(s)-v(t)$\;
$C(s,t)\leftarrow 2W\cdot R_{\mathrm{eff}}(s,t)$\;
\tcp{Hitting times to a fixed target $t$: one solve}
Solve $\Lap\,\phi=\mathbf{d}-2W\,\mathbf{1}_t$ with $\phi_t=0$ \tcp*{$\mathbf{d}$ = conductance/degree vector}
\For{each $a\in V$}{$H(a,t)\leftarrow \phi_a$\;}
$H(s,t)\leftarrow \phi_s$;\quad $H(t,s)\leftarrow C(s,t)-H(s,t)$\;
\Return $R_{\mathrm{eff}}(s,t),\,C(s,t),\,\phi$\;
\end{algorithm}

\begin{whybox}
Why does solving $\Lap\phi=\mathbf{d}-2W\mathbf{1}_t$ give \emph{all} hitting times to $t$ at once? Because the derivation in Theorem~\ref{ch35:thm-commute} showed $H(\cdot,t)$ is precisely the voltage of the ``inject-degree-everywhere, drain-at-$t$'' experiment. One linear system, $n$ hitting times. For \emph{all-pairs} commute times you need the full pseudoinverse $\Lap^+$ (then $C(s,t)=2m(\Lap^+_{ss}+\Lap^+_{tt}-2\Lap^+_{st})$), a single $\Oh(n^3)$ inversion or $n$ solves.
\end{whybox}

\begin{complexitybox}
\textbf{Cost.} Each Laplacian solve is $\Oh(n^3)$ by dense Gaussian elimination, $\Oh(m^{1.5})$-ish by conjugate gradient with a good preconditioner, or $\tilde\Oh(m)=\Oh(m\log^{c}n\,\log(1/\eps))$ by the Spielman--Teng / Koutis--Miller--Peng near-linear solvers. A single resistance/commute query: \emph{two} solves, $\tilde\Oh(m)$. All-pairs: one pseudoinverse, $\Oh(n^3)$, or Johnson--Lindenstrauss plus $\Oh(\log n/\eps^2)$ solves to estimate \emph{all} resistances to $1\pm\eps$ in $\tilde\Oh(m/\eps^2)$ total (Spielman--Srivastava).
\end{complexitybox}

\section{Mixing time and the spectral gap}\label{ch35:mixing}

The last loose end is convergence: how fast does $p_t\to\pi$? For reversible chains this is governed by the spectrum of $\Adjm_P$, which is similar to a symmetric \emph{normalized} matrix and hence to the normalized Laplacian $\mathcal{L}=\Adjm_I-\Adjm_D^{-1/2}\Adjm\Adjm_D^{-1/2}$ of the spectral chapter.

\begin{theorem}[Spectral mixing bound]\label{ch35:thm-mixing}
Let $\Adjm_P=\Adjm_D^{-1}\Adjm$ have eigenvalues $1=\mu_1>\mu_2\ge\cdots\ge\mu_n\ge-1$, and set $\mu_\star=\max(\mu_2,\abs{\mu_n})$ (the second-largest in absolute value). Then for the lazy walk and any start,
\[
  \norm{p_t-\pi}_{\mathrm{TV}}\;\le\;\frac{1}{2}\sqrt{\frac{1-\pi_{\min}}{\pi_{\min}}}\;\mu_\star^{\,t},
  \qquad
  t_{\mathrm{mix}}(\eps)=\Oh\!\Big(\frac{1}{1-\mu_\star}\log\frac{1}{\pi_{\min}\eps}\Big).
\]
The quantity $1-\mu_2$ is the \emph{spectral gap}; a larger gap means faster mixing.
\end{theorem}

\begin{proof}[Idea]
Conjugate $\Adjm_P$ to the symmetric $\Adjm_S=\Adjm_D^{1/2}\Adjm_P\Adjm_D^{-1/2}=\Adjm_D^{-1/2}\Adjm\Adjm_D^{-1/2}=\Adjm_I-\mathcal{L}$, which shares $\Adjm_P$'s eigenvalues and has an orthonormal eigenbasis. Expand the (rescaled) start vector in that basis; the $\mu_1=1$ component is the stationary part, and every other component decays by $\mu_k$ per step. The slowest survivor decays like $\mu_\star^t$. Converting the $\ell_2$ bound to total variation contributes the $\sqrt{(1-\pi_{\min})/\pi_{\min}}$ factor. Laziness ($\Adjm_P'=\tfrac12(\Adjm_I+\Adjm_P)$) maps eigenvalues to $\tfrac12(1+\mu)\in[0,1]$, killing the negative-eigenvalue (periodicity) problem so $\mu_\star=\mu_2$.
\end{proof}

By Cheeger's inequality (the expander chapter), the gap $1-\mu_2$ is sandwiched by the conductance $\Phi$: $\tfrac{\Phi^2}{2}\le 1-\mu_2\le 2\Phi$. So a graph mixes fast iff it has no sparse cut iff it is an expander. Expanders, random walks, and the Laplacian gap are three names for one phenomenon, which is why the random walk is the algorithmic heart of sampling, undirected-connectivity in logspace, and PageRank (next chapter).

\begin{labbox}
Estimate the mixing time of the running graph numerically. Build $\Adjm_P$ from Example~\ref{ch35:ex-running}, lazify it, and iterate $p_{t+1}=p_t\Adjm_P'$ from $p_0=\mathbf{1}_5$ (start at the leaf). Watch $p_t$ crawl toward $\pi=(.2,.2,.3,.2,.1)$. Plot $\norm{p_t-\pi}_1$ on a log axis; the slope is $\log\mu_\star$. Then compute $\mu_2$ of $\Adjm_P$ directly (read it off the $5\times5$ matrix or get it numerically) and check the slope matches.
\end{labbox}

\section{The full picture on one network}

It is worth seeing every quantity for the running graph side by side, computed two ways and agreeing.

\begin{center}
\begin{tabular}{lccc}
\toprule
Quantity & By hand (probability) & By resistance & Value\\
\midrule
$\pi$ & first-step / fixed point & $\deg/2m$ & $(.2,.2,.3,.2,.1)$\\
$H(3,5)$ & Example~\ref{ch35:ex-hitting} & --- & $16$\\
$H(5,3)$ & $C-H(3,5)$ & $2m R_{\mathrm{eff}}-16$ & $4$\\
$R_{\mathrm{eff}}(3,5)$ & --- & series $1{+}1$ & $2$\\
$C(3,5)$ & $H(3,5){+}H(5,3)$ & $2m\,R_{\mathrm{eff}}$ & $20$\\
$R_{\mathrm{eff}}(1,2)$ & --- & parallel $1\,\|\,2$ & $2/3$\\
$\sum_e R_{\mathrm{eff}}(e)$ & --- & Foster & $4=n{-}1$\\
\bottomrule
\end{tabular}
\end{center}

\begin{figure}[H]
\centering
\begin{tikzpicture}[font=\small,thick,>=Stealth]
  \node[draw,rounded corners,fill=green!6,minimum width=3.4cm,minimum height=1cm] (rw) at (0,2) {Random walk $\Adjm_P=\Adjm_D^{-1}\Adjm$};
  \node[draw,rounded corners,fill=blue!6,minimum width=3.4cm,minimum height=1cm] (el) at (7,2) {Resistor network $c_{ij}$};
  \node[draw,rounded corners,fill=gray!8,minimum width=3.4cm,minimum height=1cm] (la) at (3.5,0) {Laplacian $\Lap,\ \Lap^+$};
  \draw[<->] (rw) -- node[above,font=\scriptsize]{$C=2mR_{\mathrm{eff}}$} (el);
  \draw[->] (rw) -- node[left,font=\scriptsize,xshift=-2pt]{$\Lap\phi=\mathbf{d}-2m\mathbf{1}_t$} (la);
  \draw[->] (el) -- node[right,font=\scriptsize,xshift=2pt]{$\Lap v=\mathbf{1}_s-\mathbf{1}_t$} (la);
  \node[font=\scriptsize] at (3.5,-1) {hitting times, resistances, Foster, mixing $\Rightarrow$ all are Laplacian solves};
\end{tikzpicture}
\caption{The three faces of one object. Walk, circuit, and matrix are interconvertible; the commute--resistance identity is the top edge, and both sides reduce to Laplacian linear systems.}
\label{ch35:fig-triangle}
\end{figure}

\begin{historybox}
The electrical analogy is old: Kirchhoff (1847) related spanning trees to currents, and Lord Kelvin and Maxwell formalized energy minimization. The probabilistic dictionary was sharpened by Doyle and Snell's gem \emph{Random Walks and Electric Networks} (1984), which made ``voltage $=$ hitting probability'' undergraduate-accessible. The commute-time identity $C=2mR_{\mathrm{eff}}$ is due to Chandra, Raghavan, Ruzzo, Smolensky, and Tiwari (1989), who used it to bound space-bounded computation. Foster's theorem dates to 1949. The modern explosion --- Spielman--Srivastava sparsification (2008), near-linear Laplacian solvers (Spielman--Teng 2004; Koutis--Miller--Peng), and electrical max-flow (Christiano et al., 2011) --- turned resistances from a pretty analogy into the fastest known algorithms for flows and cuts, the subject of the solvers chapter.
\end{historybox}

\begin{applicationbox}
Effective resistance and commute time are working tools, not curiosities. \emph{Clustering and ML}: commute-time distance is a robust similarity that respects density (used in spectral clustering and manifold learning). \emph{Sparsification}: resistance-based sampling builds spectral sparsifiers in near-linear time. \emph{Recommendation / link prediction}: hitting and commute times rank candidate edges in networks. \emph{Sensor networks and chip design}: literal resistances. \emph{Theory}: the $\mathbf{RL}\subseteq\mathbf{L}$ line of work and Reingold's $\mathbf{SL}=\mathbf{L}$ undirected-connectivity algorithm are random-walk arguments at heart.
\end{applicationbox}

\section*{Exercises}
\addcontentsline{toc}{section}{Exercises}

\begin{exercise}[Stationary by hand]
Take the path graph $P_4$ on vertices $1\!-\!2\!-\!3\!-\!4$. Write $\Adjm_P$, find $\pi$ from $\deg/2m$, and verify $\pi\Adjm_P=\pi$ directly. Why does the (non-lazy) walk fail to converge from $X_0=1$? Identify the period.
\end{exercise}

\begin{exercise}[Hitting times on a cycle]
For the $n$-cycle $C_n$, prove by symmetry and first-step analysis that $H(i,j)=d(n-d)$ where $d$ is the graph distance from $i$ to $j$. Conclude $C(i,j)=2H$ when $i,j$ are antipodal, and check it against $2m\,R_{\mathrm{eff}}$ with $m=n$ and $R_{\mathrm{eff}}=\frac{d(n-d)}{n}$ (two arcs in parallel).
\end{exercise}

\begin{exercise}[The lollipop is worst]
The lollipop $L_n$ is a clique $K_{n/2}$ with a path of length $n/2$ attached at one clique vertex $u$, ending at $w$. Show $H(w,u)=\Th(n^2)$ (path) but $H(u,w)=\Th(n^3)$ (escape the clique, then traverse). Conclude the $\Th(n^3)$ worst-case hitting time and that hitting time can be wildly asymmetric.
\end{exercise}

\begin{exercise}[Resistance is a metric]
Using $R_{\mathrm{eff}}(s,t)=\chi^\top\Lap^+\chi$ with $\chi=\mathbf{1}_s-\mathbf{1}_t$, prove $\sqrt{R_{\mathrm{eff}}}$ is a metric: show $R_{\mathrm{eff}}(s,t)=\norm{(\Lap^+)^{1/2}\mathbf{1}_s-(\Lap^+)^{1/2}\mathbf{1}_t}^2$ and apply the Euclidean triangle inequality. Where does connectivity get used?
\end{exercise}

\begin{exercise}[Rayleigh monotonicity]
Prove that adding an edge to a graph cannot increase $R_{\mathrm{eff}}(s,t)$ for any $s,t$, and that increasing any edge's resistance cannot decrease $R_{\mathrm{eff}}(s,t)$. (Use Thomson's principle: a minimizer over a larger feasible set, or with smaller costs, can only do better.) Deduce that for a spanning subgraph $H\subseteq\Gph$, $R^{H}_{\mathrm{eff}}(s,t)\ge R^{\Gph}_{\mathrm{eff}}(s,t)$.
\end{exercise}

\begin{exercise}[Foster, twice]
Verify Foster's theorem on (a) the complete graph $K_4$ and (b) the star $K_{1,3}$. For $K_n$, show every edge has $R_{\mathrm{eff}}=2/n$ and re-derive Foster from $\binom n2\cdot\frac2n=n-1$.
\end{exercise}

\begin{exercise}[Commute time from the pseudoinverse]
Compute $\Lap^+$ for the triangle $K_3$ explicitly (it is circulant). Read off $C(1,2)=2m(\Lap^+_{11}+\Lap^+_{22}-2\Lap^+_{12})$ and check against $2m\,R_{\mathrm{eff}}(1,2)$ with $R_{\mathrm{eff}}=2/3$, $m=3$. You should get $C=4$.
\end{exercise}

\begin{exercise}[Harmonic = hitting probability, coded]
Implement the Dirichlet solve: given a graph and two vertices $s,t$, set $v(s)=1,v(t)=0$ and solve $(\Lap v)(a)=0$ at interior $a$. Then estimate the same hitting probabilities by Monte Carlo simulation of the walk. Plot solver vs.\ simulation; they should agree within sampling noise. (This is the cheapest empirical proof of Theorem~\ref{ch35:thm-dirichlet}.)
\end{exercise}

\begin{exercise}[Spectral gap and a barbell]
The barbell graph is two cliques $K_{n/2}$ joined by a single edge. Argue informally that $\mu_2$ is very close to $1$ (a near-disconnected cut), so mixing is slow and the commute time across the bridge is $\Th(n^3)$. Relate this to Cheeger: the bottleneck edge gives tiny conductance $\Phi=\Th(1/n^2)$.
\end{exercise}

\begin{exercise}[Weighted walks]
Generalize Theorem~\ref{ch35:thm-commute} to a weighted graph with conductances $c_{ij}$: define the walk with $\Pr[i\to j]=c_{ij}/c_i$, set $2W=\sum_i c_i$, and prove $C(s,t)=2W\,R_{\mathrm{eff}}(s,t)$ by repeating the two-experiment subtraction. Identify exactly where reversibility (detailed balance) is used.
\end{exercise}

\begin{exercise}[Cover time of the path]
Show the cover time of the path $P_n$ is $\Th(n^2)$ (the walk must reach both ends; the far end is $\Th(n^2)$ away). Then use the Matthews bound with $H_{\max}=\Th(n^2)$ to get the loose upper bound $\Oh(n^2\log n)$, and explain why for the path the $\log n$ is \emph{not} tight, whereas for the complete graph cover time $\Th(n\log n)$ it is.
\end{exercise}
